\subsection{The Peter-Weyl theorem}

We recall that we denote by \(\Theta(\mathrm{G})\) the space of matrix coefficients of finite-dimensional unitary representations of G. The space \(\Theta(\mathrm{G})\) is a unital subalgebra of \(\mathcal{C}(\mathrm{G})\) stable under complex conjugation (prop. 5 of V, p. 386).

For \(\pi \in \widehat{\mathrm{G}}\) , let \(\varrho_{\mathrm{G}}(\pi)\) be the subspace of \(\mathcal{C}(\mathrm{G})\) generated by the matrix coefficients of \(\pi\) . We identify it with a subspace of \(\mathrm{L}^2 (\mathrm{G})\) .

The space \(\Theta(G)\) coincides with the sum of the spaces \(\varrho_{G}(\pi)\) for \(\pi\in\widehat{G}\) . Indeed, every element of \(\Theta(G)\) is a sum of matrix coefficients of irreducible representations of G, since the finite-dimensional representations of G are semisimple by Prop. 1 of V, p. 457). Moreover, this sum is direct since, by Cor. 2 of V, p. 458, the spaces \(\varrho_{G}(\pi)\) are pairwise orthogonal.

Proposition 4. — The space \(\Theta(\mathrm{G})\) is dense in \(\mathcal{C}(\mathrm{G})\) and in \(\mathrm{L}^2(\mathrm{G})\) .

The unital subalgebra \(\Theta(\mathrm{G})\) of \(\mathcal{C}(\mathrm{G})\) is stable under conjugation. It coincides with the subalgebra \(\Upsilon(\mathrm{G})\) by prop. 2 of V, p. 457, hence it separates the points of G (cor. of th. 4 of V, p. 454). Consequently, it is dense in \(\mathcal{C}(\mathrm{G})\) by TG, X, p. 40, cor. 2, and a fortiori, it is dense in the space \(\mathrm{L}^2(\mathrm{G})\) (cf. INT, IV, p. 155, §4, no. 7, prop. 13).

Recall that the biregular unitary representation of G is the representation \(\varrho_{G}\) of G × G into \(\mathrm{L}^{2}(\mathrm{G})\) such that \((g_{1}, g_{2}) \mapsto \gamma_{\mathrm{G}}(g_{1}) \delta_{\mathrm{G}}(g_{2})\) .

PROPOSITION 5. --- Let \(\pi \in \widehat{\mathbf{G}}\) . The space \(\varrho_{\mathrm{G}}(\pi)\) is a subrepresentation of \(\varrho_{\mathrm{G}}\) which is isomorphic to \(\overline{\pi} \boxtimes \pi\) . In particular, it is an irreducible representation of \(\mathrm{G} \times \mathrm{G}\) .

By prop. 7 of V, p. 422 (applied with \(Z = \{e\}\) ), the space \(\varrho_{G}(\pi)\) is a subrepresentation of \(\varrho_{G}\) and the linear map from \(\overline{\mathrm{E}}_{\pi} \otimes \mathrm{E}_{\pi}\) to \(L^2(G)\) that sends \(x \otimes y\) to the matrix coefficient \(g \mapsto \langle x | \pi(g)y \rangle\) is a \((G \times G)\)-isomorphism from \(\overline{\pi} \boxtimes \pi\) onto \(\varrho_{G}(\pi)\) . The subrepresentation \(\varrho_{G}(\pi)\) is therefore irreducible (prop. 3 of V, p. 461).

THEOREM 1 (Peter--Weyl). --- Let G be a compact topological group. The biregular representation \(\varrho_{\mathrm{G}}\) of G is the Hilbert sum of the subrepresentations \(\varrho_{\mathrm{G}}(\pi)\) for \(\pi \in \widehat{\mathrm{G}}\) .

The spaces \(\varrho_{\mathrm{G}}(\pi)\) are pairwise orthogonal; these are irreducible subrepresentations of the biregular representation of G (prop. 5) which are pairwise non-isomorphic (prop. 3 of V, p. 461).

The theorem then follows from the fact that the sum \(\Theta(\mathrm{G})\) of the spaces \(\varrho_{\mathrm{G}}(\pi)\) for \(\pi\) running through \(\widehat{\mathrm{G}}\) is dense in \(\mathrm{L}^2(\mathrm{G})\) (prop. 4).

COROLLARY 1. --- For every \(\pi\in\widehat{G}\) , the subrepresentation \(\varrho_{\mathrm{G}}(\pi)\) is the \((\overline{\pi}\boxtimes\pi)\)-isotypic component of \(\varrho_{\mathrm{G}}\) .

Let F be the \((\overline{\pi} \boxtimes \pi)\)-isotypic component of \(\varrho_{G}\) . The space F contains \(\varrho_{\mathrm{G}}(\pi)\) (prop. 5). Moreover, for every \(\tau \in \widehat{G}\) different from \(\pi\) , the intersection of F and \(\varrho_{\mathrm{G}}(\tau)\) is zero (prop. 3 of V, p. 461). From this we deduce that \(\mathrm{F} = \varrho_{\mathrm{G}}(\pi)\) by applying the theorem.

COROLLARY 2. --- Let \(\pi_1\) and \(\pi_2\) be non-isomorphic irreducible representations of G. The \((\overline{\pi}_1 \boxtimes \pi_2)\)-isotypic component of \(\varrho_G\) is zero.

COROLLARY 3. --- The right (resp. left) regular representation of G is isomorphic to the Hilbert sum

\[
\bigoplus_ {\pi \in \widehat {\mathrm{G}}} \pi^ {\dim (\pi)}.
\]

By the theorem and proposition 5, the restriction of the biregular representation of G to the subgroup \(\mathrm{H} = \{e\} \times \mathrm{G}\) is isomorphic to the Hilbert sum of the restrictions to H of the representations \(\overline{\pi} \boxtimes \pi\) for \(\pi \in \widehat{\mathrm{G}}\) . These are isomorphic to the direct sum of \(\dim(\pi)\) copies of \(\pi\) (lemma 8 of V, p. 384). This implies the result for the right regular representation, and the case of the left regular representation is similar.

COROLLARY 4. --- Let \(\pi\in\widehat{G}\) . The \(\pi\)-isotypic component of \(\delta_{G}\) (resp. the component \(\overline{\pi}\)-isotypic of \(\gamma_{G}\) ) is equal to \(\varrho_{G}(\pi)\) .

The argument is similar to that of the preceding corollary, by considering \(\varrho_{\mathrm{G}}(\pi)\) as a subrepresentation of \(\delta_{\mathrm{G}}\) (resp. of \(\gamma_{\mathrm{G}}\) ).

Remarks. --- 1) If G is finite, these statements correspond to the results of A, VIII, p. 398, remark.

\begin{enumerate}
\def\labelenumi{\arabic{enumi})}
\setcounter{enumi}{1}
\item
  Suppose that G is commutative. The set \(\widehat{G}\) identifies with the dual group of G (V, p. 393, remark). For every \(\chi \in \widehat{G}\) , the space \(\varrho_{\mathrm{G}}(\chi)\) is the 1-dimensional subspace of \(\mathrm{L}^2 (\mathrm{G})\) generated by \(\chi\) . Theorem 1 for G is therefore equivalent to the corollary of Theorem 1 of II, p. 215.
\item
  If there exists an integer \(n \geqslant 0\) such that G is a compact subgroup of \(\mathbf{GL}(\mathbf{C}^{n})\) , the identity representation of G in \(\mathbf{GL}(\mathbf{C}^{n})\) is enough to separate the points of G, and one can then prove Proposition 4 directly, followed by the Peter–Weyl theorem, without invoking the Gelfand–Raikov theorem.
\item
  The Peter–Weyl theorem implies in particular that the natural continuous homomorphism from G into \(\prod_{\pi\in\widehat{G}}\mathbf{U}(\mathrm{E}_{\pi})\) is injective.
\end{enumerate}

